\section{関連研究}
\label{sec:related_work}

関連研究は二本柱で整理する。
第一の柱はスポーツフィールド・コートの登録と校正であり、
画像から何を取り出し、既知のフィールド形状へどう対応付け、何を出力するかの系譜である。
第二の柱はニューラル三次元再構成・合成データの下流利用であり、
どのような学習データをどう作るかの系譜である。
本研究はどちらの柱も置き換えず、
\emph{学習データの視点分布}と、
その分布を作るための\emph{メートル系アライメント}に焦点を置く。

\subsection{スポーツフィールドの表現と校正}

スポーツ映像のカメラ校正は、
画像から抽出したマーキング表現を既知のフィールド形状へ対応付ける問題として発展してきた。
深い構造モデルによるフィールド位置推定\cite{homayounfar2017sportsfield}は、
検出と構造化されたマッチングを組み合わせる初期の代表例である。
合成データによる校正\cite{chen2019sports}は、
既知形状を多数のカメラから投影した線画像と姿勢辞書を使って放送画像を校正する。
ここで合成されるのは実会場の写実的な RGB ではなく、
主としてフィールド形状を表す線テンプレートである点が、本研究との差分になる。

領域表現と学習型の補正を組み合わせる方法\cite{sha2020endtoend}は、
意味的な領域分割と相対的なホモグラフィ補正を結び付ける。
学習した誤差モデルを通して校正を最適化する方法\cite{jiang2020optimizing}や、
格子状のキーポイントと距離マップを共有ネットワークで扱う登録\cite{nie2021robust}は、
複数の幾何教師を共有表現へ与える系統である。
テニスを含む評価を行った研究\cite{nie2021robust}は、
良好な結果を小さなカメラ移動やコート全体が写る場面の多さと関連付けており、
視点条件が性能を左右することを示唆する。

キーポイントと線の対応付けを明示的に扱う系統も存在する。
意味付きフィールド要素と既知三次元モデルの投影距離でカメラを最適化する TVCalib\cite{theiner2023tvcalib}、
点と線の再投影誤差で校正を補正する PnLCalib\cite{gutierrez2026pnlcalib}、
任意のキーポイント分割を用いる一手の登録\cite{jacquelin2022oneshot}、
キーポイント条件付きのラベル設計\cite{chu2022kpsfr}がこれにあたる。
本研究の検出器は検出した線を入力して姿勢を解く直列構成ではなく、
共有特徴から密な予測と姿勢を並列に回帰する。
したがってこれらのソルバを本研究が内蔵しているとは主張せず、
\emph{点と線の幾何を一つの仮説へまとめる}という問題意識と、
学習データの視点分布という差分のみを主張する。

古典的なコートモデル当てはめ\cite{farin2004courtmodels}は、
規制寸法のテンプレートを画像へ合わせるという本研究のアライメントと近い発想を持つ。
本研究のアライメントは画像平面ではなく、
再構成シーンのメートル系座標でテンプレートを当てはめる点が異なる。
非放送・実用撮影を扱う研究も存在する。
移動する撮影機器を含むスポーツフィールドのカメラ姿勢推定\cite{citraro2020realtime}や、
アマチュア撮影のテニス映像を対象とした線検出\cite{agrawal2024amateur}は、
放送主カメラ以外の条件が実在の課題であることを示す。
したがって「従来研究は放送映像のみを扱ってきた」とは記述しない。
本研究の差分は、実写収録に依存する視点分布を再構成ベースで拡張し、
その効果を画像数や外観の変更と切り分けて評価する設計にある。

評価手続き自体も研究対象である。
コート平面のホモグラフィだけでは立体構造やレンズモデルの能力を評価しきれないこと、
既知三次元物体の再投影を用いる校正非依存の評価が必要であることが指摘されている\cite{magera2024procc}。
中央視点の追加が分布外の挙動をどう変えるかを調べる研究\cite{magera2025centralviews}は、
追加画像にカメラ真値が無い場合の評価の限界も示している。
放送カメラの追跡を主題とする研究\cite{magera2025broadtrack}は、
校正を一時点の推定ではなく時間方向の一貫性として扱う。
本研究はこれらの指摘を踏まえ、
検出精度と幾何登録の精度を同じ指標へ混ぜず、
実写 held-out validation と合成 trajectory-group test を別の列として報告する。

\subsection{3DGS・NeRF と合成データの下流利用}

3DGS は異方性ガウスで放射輝度場を表現し、
リアルタイムの新規視点合成を可能にする\cite{kerbl2023gaussian}。
ニューラル再構成を検出器や姿勢推定器の学習データ生成に使う研究は、
単なる見た目の合成ではなく、生成されたデータが下流タスクへ与える影響を測る点に特徴がある。
NeRF からのデータ生成と下流タスクを微分可能な閉ループに置く Neural-Sim\cite{ge2022neuralsim}は、
「どの視点を生成するか」という分布設計の問題を明示する。
NeRF 合成画像で位置推定器を学習させる LENS\cite{moreau2021lens}や、
3DGS と物理シミュレーションから 6DoF 物体姿勢のデータセットを作る PEGASUS\cite{meyer2024pegasus}は、
再構成と多視点 RGB・幾何教師の組み合わせが既に有用であることを示している。
近年は 3DGS を 3D 物体検出のためのデータ生成器として評価する研究も現れている\cite{zanjani2025gsgenerator}。
物体を切り出して背景画像へ合成する系統\cite{vanherle2025cutandsplat}を含め、
「3DGS で下流用データを生成する」こと自体を本研究は新規性として主張しない。

手続き型の合成データ生成も同様に確立している。
Kubric\cite{greff2022kubric}は大規模で再現可能なデータセット生成を目指し、
BlenderProc\cite{denninger2019blenderproc}はその実装基盤として広く使われている。
スポーツのフィールド検出に限っても、
シミュレータ由来の RGB で検出器を学習させる研究\cite{qin2025soccersynth}が存在する。
この系統では合成 RGB による検出改善自体が既に報告されているため、
本研究の差分は「合成データでフィールド検出を改善する」ことではなく、
\emph{実会場の再構成}を出発点にすること、
および\emph{撮影位置・方向の分布}を制御して
外観変更や画像数増加と切り分ける設計を置くことにある。

\subsection{本研究の位置づけと引用上の注意}

以上の整理から、本研究は次の三点で先行研究と区別される。

\begin{enumerate}[leftmargin=1.6em]
  \item 学習データはシミュレータの仮想会場ではなく、
  対象ドメインの実映像から再構成したシーンと、
  そこで一度だけ検証されたメートル系コートアライメントから作る。
  \item 生成画像は背景の差し替えではなく、
  復元カメラの凸包で制限した軌道からのレンダ結果である。
  したがって外観は撮影会場に固定され、
  可変なのは撮影位置・方向・可視範囲と対象コートである。
  \item 単一のシーン--コート変換から、
  キーポイント、領域、線、線の意味分類、カメラ姿勢を同じ画像フレームで生成する。
  教師どうしの不整合はデータ生成の段階で生じない。
\end{enumerate}

引用に際しては次の版管理上の注意を置く。
PnLCalib\cite{gutierrez2026pnlcalib}は CVPRW 2024 の発表
\emph{No Bells, Just Whistles} の拡張版であり、
同一 arXiv 系列の初版と拡張版を独立した二つの発明として数えない。
TVCalib\cite{theiner2023tvcalib}の比較には主カメラ中心の評価部分集合が使われており、
本研究で参照する場合も対象集合を明示する。
SoccerSynth\cite{qin2025soccersynth}の条件では
複数カメラと地面色が同時に変わるため、
本研究では視点のみの寄与を切り分ける条件を別に用意する。
公開実装である \BaselineName\cite{tcdrepo} は
本研究が実写 source と比較相手として実際に用いるもので、
付随する論文の有無や再配布条件は付録~\ref{app:reproducibility} に記す。
